% !TEX program = pdflatex
\documentclass[10pt]{article}
\usepackage[utf8]{inputenc}\usepackage[T1]{fontenc}\usepackage[english]{babel}
\usepackage{lmodern,microtype}\usepackage[margin=.72in]{geometry}
\usepackage{amsmath,amssymb,amsthm,mathtools,booktabs,array,xcolor,enumitem}
\usepackage[most]{tcolorbox}\usepackage{fancyhdr}\usepackage[hidelinks]{hyperref}
\definecolor{CourseBlue}{HTML}{1A3C68}\definecolor{CourseGreen}{HTML}{216E39}
\definecolor{CourseGold}{HTML}{B8860B}\definecolor{CourseRed}{HTML}{9A2C2C}
\definecolor{SoftBlue}{HTML}{F2F6FA}\definecolor{SoftGreen}{HTML}{F1F8F3}
\definecolor{SoftGold}{HTML}{FFF8E8}\definecolor{SoftRed}{HTML}{FFF3F3}
\setlength{\parindent}{0pt}\setlength{\parskip}{.32em}
\setlist[itemize]{leftmargin=1.5em,itemsep=.12em,topsep=.2em}
\pagestyle{fancy}\fancyhf{}\lhead{Stochastic Calculus and Option Pricing}
\rhead{Lecture 13 Notes}\cfoot{\thepage}\renewcommand{\headrulewidth}{.3pt}
\newcommand{\E}{\mathbb E}\newcommand{\Q}{\mathbb Q}\newcommand{\cF}{\mathcal F}
\newcommand{\dd}{\,\mathrm d}
\theoremstyle{definition}\newtheorem{definition}{Definition}[section]
\newtcolorbox{boardexample}[1]{enhanced,breakable,colback=SoftBlue,colframe=CourseBlue,boxrule=.55pt,arc=1mm,title={Example --- #1},fonttitle=\bfseries}
\newtcolorbox{studentexercise}[1]{enhanced,breakable,colback=SoftGold,colframe=CourseGold,boxrule=.55pt,arc=1mm,title={Exercise --- #1},fonttitle=\bfseries}
\newtcolorbox{shortsolution}[1][]{enhanced,breakable,colback=SoftGreen,colframe=CourseGreen,boxrule=.55pt,arc=1mm,title={Short solution#1},fonttitle=\bfseries}
\newtcolorbox{keyidea}[1]{enhanced,breakable,colback=white,colframe=CourseGreen,boxrule=.7pt,arc=1mm,title={#1},fonttitle=\bfseries}
\newtcolorbox{warningbox}[1]{enhanced,breakable,colback=SoftRed,colframe=CourseRed,boxrule=.65pt,arc=1mm,title={#1},fonttitle=\bfseries}
\title{\textbf{Pricing, Hedging, and Model Risk}\\\large Lecture 13 Notes --- Vanilla, Path Dependence, and Stochastic Rates}
\author{}\date{}
\begin{document}\maketitle\vspace{-1.4em}
\begin{keyidea}{Goals and logical structure}
Pricing is only the beginning of a derivatives workflow:
\[
\text{contract}\longrightarrow\text{state variables}\longrightarrow
\text{model and price}\longrightarrow\text{hedge}\longrightarrow\text{model-risk audit}.
\]
This lecture compares three situations: a vanilla option, a lookback option,
and a claim exposed to stochastic interest rates.
\end{keyidea}

\section{A vanilla option: the hedge depends on the model}
For a European payoff \(g(S_T)\), Black--Scholes assumes
\[
\frac{\dd S_t}{S_t}=r\dd t+\sigma\dd W_t^\Q.
\]
The price is
\[
v(t,s)=\E^\Q_{t,s}[e^{-r(T-t)}g(S_T)],
\]
and the delta hedge holds \(\Delta_t=v_s(t,S_t)\) shares. In this
one-factor complete model, continuous rebalancing replicates the claim.

\begin{boardexample}{local volatility changes delta}
Under local volatility,
\[
\frac{\dd S_t}{S_t}=r\dd t+\sigma_{\rm loc}(t,S_t)\dd W_t^\Q,
\]
and
\[
v_t+rs v_s+\tfrac12\sigma_{\rm loc}(t,s)^2s^2v_{ss}-rv=0.
\]
The formal hedge remains \(\Delta=v_s\), but its value and delta depend on
the calibrated surface. A hedge is therefore a model-dependent object.
\end{boardexample}

\begin{warningbox}{A good vanilla fit is not a guarantee of a good hedge}
If the market follows stochastic volatility while the desk uses a
local-volatility delta, the option price may initially match market vanillas
but the subsequent smile movement can generate systematic hedging P\&L.
Calibration error, discrete rebalancing, transaction costs, and parameter
instability add further sources of error.
\end{warningbox}

\begin{studentexercise}{price versus hedge}
Two models match today's vanilla call price but imply different deltas.
Can their future hedging P\&L differ?
\end{studentexercise}
\begin{shortsolution}
Yes. Equality of one initial price does not imply equality of future
sensitivities or of model-implied smile dynamics. Hedging P\&L depends on
the path followed by market risk factors and on the hedge rule.
\end{shortsolution}

\section{A lookback option: the path becomes state}
\subsection{The additional state variable}
A fixed-strike lookback put may pay
\[
M_T-S_T,\qquad M_t=\max_{0\leq u\leq t}S_u.
\]
Knowing \(S_t\) alone is insufficient: two paths with equal current spot but
different running maxima have different future payoffs. The Markov state is
\((S_t,M_t)\).

\begin{definition}[State augmentation]
When a payoff depends on a path statistic, include the statistic as a state
variable. Pricing then becomes Markovian in the enlarged state, not
necessarily in the stock price alone.
\end{definition}

Away from a new maximum, a local-volatility lookback value \(u(t,s,m)\)
satisfies
\[
u_t+rsu_s+\tfrac12\sigma_{\rm loc}(t,s)^2s^2u_{ss}-ru=0,
\qquad 0<s<m,
\]
with terminal condition \(u(T,s,m)=m-s\). The diagonal \(s=m\) requires a
boundary condition describing the update of the running maximum.

\begin{keyidea}{Path dependence has a computational cost}
The new state adds a grid dimension and makes simulation bookkeeping
essential. In Monte Carlo, each path must carry its own running maximum;
in a PDE, the diagonal boundary must be treated consistently.
\end{keyidea}

\begin{studentexercise}{minimal state}
For an arithmetic Asian call with payoff
\((T^{-1}\int_0^T S_u\dd u-K)^+\), what extra state restores Markovianity?
\end{studentexercise}
\begin{shortsolution}
Use the running integral \(I_t=\int_0^tS_u\dd u\), so the state is
\((t,S_t,I_t)\). The accumulated average alone does not replace the current
stock because future accumulation still depends on \(S_t\).
\end{shortsolution}

\section{Stochastic rates add a hedging factor}
If the short rate is random, discounting is random:
\[
P(t,T)=\E^\Q\left[\exp\left(-\int_t^T r_u\dd u\right)\middle|\cF_t\right].
\]
For a stock and a short-rate factor,
\[
\frac{\dd S_t}{S_t}=(r_t-q)\dd t+\sigma_S\dd W_t^S,\qquad
\dd r_t=a(t,r_t)\dd t+\sigma_r(t,r_t)\dd W_t^r.
\]
The value of a European claim is generally \(v(t,S_t,r_t)\). Its hedge
exposure includes both \(v_s\) and \(v_r\).

\begin{boardexample}{what becomes tradable?}
Stock delta manages the stock shock. A zero-coupon bond or an interest-rate
instrument is needed to manage the rate shock. If sufficiently many
independent risk factors cannot be spanned by traded instruments, the
market is incomplete and perfect replication is unavailable.
\end{boardexample}

\begin{warningbox}{Model risk and market incompleteness are different}
Model risk is error from a misspecified or unstable model. Incompleteness is
a structural fact: available traded assets cannot span all shocks. Both can
produce residual P\&L, but their remedies differ.
\end{warningbox}

\section{Hedging P\&L as a diagnostic}
For a short option position, a simple discrete hedging P\&L is
\[
\Pi_T=V_0+\sum_i\Delta_{t_i}(S_{t_{i+1}}-S_{t_i})
-\text{option payoff}-\text{funding and costs}.
\]
Analyse not only its mean, but also dispersion, tail losses, dependence on
spot moves, volatility moves, and rate moves. A model can have small average
error while producing unacceptable tail risk.

\begin{keyidea}{Practical workflow}
\begin{enumerate}
\item Identify every payoff-relevant state variable.
\item Calibrate the model to liquid instruments that support hedging.
\item State the hedge instruments and rebalance convention.
\item Backtest and stress-test the distribution of hedging P\&L.
\item Report residual risk rather than claiming a perfect hedge by default.
\end{enumerate}
\end{keyidea}

\begin{studentexercise}{final diagnosis}
A long-dated lookback option is priced with constant volatility and
deterministic rates. Name two material modeling omissions to investigate.
\end{studentexercise}
\begin{shortsolution}
The path-dependent payoff can be very sensitive to volatility-surface and
smile dynamics, so constant volatility is a key omission. Long maturity also
makes stochastic-rate and stock-rate correlation risk potentially material.
The correct choice depends on available market instruments and the P\&L
impact of these risks.
\end{shortsolution}
\end{document}
